% Part: sets-functions-relations
% Chapter: arithmetization
% Section: checking-details
%
\documentclass[../../../include/open-logic-section]{subfiles}

\begin{document}

\olfileid[fr]{sfr}{arith}{check}
\olsection{Anneaux et corps ordonnés}

Tout au long de ce chapitre, nous avons affirmé que certaines définitions
donnaient aux objets construits les propriétés « attendues ». Dans cette
annexe technique, nous allons préciser ce que nous entendons par là et
montrer, ou esquisser comment montrer, qu'elles ont bien cet effet.

Dans la \olref[int]{sec}, nous avons défini l'addition et la multiplication
sur $\Int$. Nous voulons montrer que ces opérations, telles que nous
les avons définies, munissent $\Int$ de la structure attendue.
Il s'agit, plus précisément, d'une structure d'anneau commutatif.

\Needspace{16\baselineskip}\begin{defn}
Un \emph{anneau commutatif} est un ensemble $S$ muni de deux éléments
distingués $0$ et $1$ et de deux opérations $+$ et $\times$ de
$S\times S$ dans $S$, qui satisfont les huit formules suivantes :
\begin{align*}
  \emph{Associativité}&&a + (b+ c) & = (a + b) + c \\
  && (a \times b) \times c & = a \times (b\times c)\\
  \emph{Commutativité}&&a + b &= b+ a  \\
  && a \times b&= b\times a\\
  \emph{Éléments neutres}&&a + 0 &= a \\
  && a \times 1 &= a\\
  \emph{Existence d'un opposé}&&(\exists b\in S)0&=a + b\\
  \emph{Distributivité}&&a \times (b+ c ) &= (a \times b) + (a \times c)
\end{align*}
Dans chaque formule, les variables libres sont implicitement quantifiées
universellement sur $S$ ; le $b$ de la formule sur l'opposé est déjà lié
par son quantificateur existentiel. Les éléments $0$ et~$1$ ne sont pas
nécessairement les nombres naturels qui portent ces noms.
\end{defn}

Pour vérifier que les entiers relatifs forment un anneau commutatif, il
suffit donc de vérifier ces huit conditions. Aucune n'est difficile à
établir, mais l'ensemble des vérifications est un peu laborieux. Voici,
par exemple, comment démontrer l'\emph{associativité} de l'addition :

\begin{proof}
Fixons $i, j, k \in \Int$. Il existe donc
$a_1, b_1, a_2, b_2, a_3, b_3 \in \Nat$ tels que
$i = \equivrep{a_1, b_1}{}$, $j = \equivrep{a_2,b_2}{}$ et
$k = \equivrep{a_3, b_3}{}$. Pour alléger les notations, nous écrivons
« $\equivrep{x, y}{}$ » au lieu de
« $\equivrep{x, y}{\Intequiv}$ » dans toute cette section. On a alors :
\begin{align*}
  i + (j + k) &= \equivrep{a_1, b_1}{}+(\equivrep{a_2, b_2}{} + \equivrep{a_3, b_3}{}) \\
  &= \equivrep{a_1, b_1}{} + \equivrep{a_2+a_3, b_2+b_3}{}\\
  &= \equivrep{a_1 + (a_2 + a_3), b_1 + (b_2 + b_3)}{}\\
  &= \equivrep{(a_1 + a_2) + a_3, (b_1 + b_2) + b_3}{}\\
  &= \equivrep{a_1 + a_2, b_1 + b_2}{} + \equivrep{a_3, b_3}{}\\
  &= (\equivrep{a_1, b_1}{} + \equivrep{a_2, b_2}{}) + \equivrep{a_3, b_3}{}\\
  &= (i+j) + k
\end{align*}
en utilisant librement les propriétés de l'addition dans $\Nat$.
\end{proof}

De même, voici comment établir l'\emph{existence d'un opposé} :

\begin{proof}
Fixons $i \in \Int$ ; on a donc $i = \equivrep{a,b}{}$ pour certains
$a,b \in \Nat$. Posons $j = \equivrep{b,a}{} \in \Int$.
Les propriétés des naturels donnent $(a+b) + 0 = 0 + (a+b)$.
Par définition, $\tuple{a+b, b+a} \sim_\Int \tuple{0,0}$, d'où
$\equivrep{a+b, b+a}{} = \equivrep{0, 0}{} = 0_\Int$.
Ainsi, $i + j = \equivrep{a,b}{}+\equivrep{b,a}{}
=\equivrep{a+b, b+a}{}=\equivrep{0,0}{} = 0_\Int$.
\end{proof}

Voici enfin une démonstration de la \emph{distributivité} :

\begin{proof}
Comme ci-dessus, fixons $i = \equivrep{a_1, b_1}{}$,
$j = \equivrep{a_2,b_2}{}$ et $k = \equivrep{a_3, b_3}{}$. Alors :
\begin{align*}
  i \times (j + k)
  &= \equivrep{a_1, b_1}{} \times (\equivrep{a_2,b_2}{} + \equivrep{a_3, b_3}{})\\
  &= \equivrep{a_1, b_1}{} \times \equivrep{a_2 + a_3,b_2+b_3}{}\\
  &= \equivrep{a_1 (a_2 + a_3) + b_1 (b_2+b_3), a_1 (b_2 + b_3) + b_1 (a_2 + a_3)}{}\\
  &= \equivrep{a_1 a_2 + a_1a_3 + b_1 b_2+b_1b_3, a_1 b_2 + a_1b_3 + a_2b_1 + a_3b_1}{}\\
% Ligne inactive de la source, erronée par répétition du terme b_1b_2 :
% &= \equivrep{a_1 a_2 + b_1b_2 + a_1a_3 + b_1 b_2+b_1b_3, a_1 b_2 + a_1b_3 + a_2b_1 + a_3b_1}{}\\
  &= \equivrep{a_1a_2 + b_1b_2, a_1b_2 + a_2b_1}{} + \equivrep{a_1a_3 + b_1b_3, a_1b_3 + a_3b_1}{}\\
  &= (\equivrep{a_1, b_1}{} \times \equivrep{a_2,b_2}{}) + (\equivrep{a_1, b_1}{} \times \equivrep{a_3, b_3}{})\\
  &= (i \times j) + (i \times k)
\end{align*}
\end{proof}

Nous laissons en exercice la démonstration des cinq autres conditions.
Une fois ces vérifications faites, nous aurons montré que $\Int$ forme
un anneau commutatif : l'addition et la multiplication définies plus
haut ont bien les propriétés attendues.

\begin{prob}
Démontrer que $\Int$ est un anneau commutatif.
\end{prob}

Le travail n'est toutefois pas terminé. Outre l'addition et la
multiplication sur $\Int$, nous avons défini une relation d'ordre
$\leq$ ; il faut vérifier qu'elle aussi a les propriétés attendues.
Plus précisément, il faut montrer que $\Int$ est un anneau
\emph{ordonné}.\footnote{Rappelons, d'après la
\olref[sfr][rel][ord]{def:linearorder}, qu'un ordre total est une
relation réflexive, transitive, antisymétrique et connexe. Pour un tel
ordre, la connexité donne la \emph{trichotomie} de l'ordre strict
associé : pour tous $a$ et $b$, exactement l'une des relations
$a < b$, $a = b$, $a > b$ est vérifiée.
\emph{Précision éditoriale :} la source écrit ici
$a \leq b \lor a = b \lor a \geq b$ ; cette disjonction est vraie,
mais ses trois cas ne sont pas exclusifs.}

\begin{defn}
Un \emph{anneau ordonné} est un anneau commutatif muni d'une relation
d'ordre total $\leq$ telle que, pour tous $a,b,c\in S$ :
\begin{align*}
  a \leq b &\lif a + c \leq b + c\\
  (a \leq b \land 0 \leq c) &\lif a \times c \leq b \times c
\end{align*}
\end{defn}

\begin{prob}
Démontrer que $\Int$ est un anneau ordonné.
\end{prob}

Comme précédemment, montrer que $\Int$, tel que nous l'avons construit,
est un anneau ordonné demande des vérifications laborieuses mais
usuelles. Nous vous les laissons.

Cela règle le cas des entiers relatifs. Il faut maintenant établir des
propriétés très semblables pour les rationnels. Plus précisément, nous
devons montrer qu'avec nos définitions de $+$, $\times$ et $\leq$, les
rationnels forment un \emph{corps} ordonné :
\begin{defn}\ollabel{orderedfield}
Un \emph{corps ordonné} est un anneau ordonné $S$ dans lequel
$0\neq1$ et qui satisfait en outre :
\begin{align*}
  \emph{Existence d'un inverse}& &
  (\forall a \in S \setminus \{0\})(\exists b \in S) a\times b& = 1
\end{align*}
\end{defn}
\begin{editorial}
La condition $0\neq1$, omise dans la source, exclut l'anneau réduit à un
seul élément. Sans elle, la condition sur les éléments non nuls serait
vide et ne suffirait pas à définir un corps.
\end{editorial}

Une fois établi que $\Int$ est un anneau ordonné, il est facile, quoique
laborieux, de montrer que $\Rat$ est un corps ordonné.

\begin{prob}
Démontrer que $\Rat$ est un corps ordonné.
\end{prob}

Après les entiers relatifs et les rationnels, il ne reste que les réels.
Il faut montrer que $\Real$ est un corps ordonné \emph{complet},
c'est-à-dire un corps ordonné ayant la propriété de la borne supérieure.
Or, le \olref[cuts]{realcompleteness} établit déjà cette propriété pour
$\Real$. Il reste toutefois à effectuer les vérifications, fastidieuses,
qui montrent que $\Real$ est un corps ordonné.

Avant d'entreprendre cet exercice laborieux, il faut vérifier quelques
faits plus « immédiats ». Par exemple, nous devons nous assurer que
$\alpha + \beta$, tel que nous l'avons défini, est bien une
\emph{coupure} dès que $\alpha$ et $\beta$ sont des coupures.
En voici une démonstration :

\begin{proof}
Puisque $\alpha$ et $\beta$ sont des coupures,
$\alpha + \beta = \Setabs{p + q}{p \in \alpha \land q \in \beta}$
est une partie non vide et propre de $\Rat$. En effet, chacune des
deux coupures contient un rationnel, dont la somme appartient à
$\alpha+\beta$. Choisissons aussi $r\in\Rat\setminus\alpha$ et
$s\in\Rat\setminus\beta$. Tout $p\in\alpha$ vérifie $p<r$ et tout
$q\in\beta$ vérifie $q<s$, car les coupures sont des segments initiaux.
Ainsi $p+q<r+s$, et $r+s\notin\alpha+\beta$.
Supposons maintenant que $x < p + q$ pour certains $p \in \alpha$ et
$q \in \beta$, avec $x\in\Rat$. Alors $x - p < q$, donc
$x - p \in \beta$, et $x = p + (x - p) \in \alpha + \beta$.
La somme est donc un segment initial de $\Rat$.
Enfin, pour tout $p + q \in \alpha + \beta$, les deux coupures étant
sans maximum, il existe $p_1 \in \alpha$ et $q_1 \in \beta$ tels que
$p < p_1$ et $q < q_1$. Ainsi
$p + q < p_1 + q_1 \in \alpha + \beta$, et $\alpha+\beta$ n'a pas de maximum.
\end{proof}

Des vérifications analogues permettent de montrer que $\alpha - \beta$,
$\alpha \times \beta$ et $\alpha \div \beta$ sont des coupures, en
excluant, dans le dernier cas, celui où $\beta$ est la coupure nulle.
Là encore, nous vous laissons ces vérifications.
\begin{editorial}
La source n'active pas ses définitions de la soustraction et de la
division des coupures. Pour donner un sens précis à cet exercice,
on pose $\alpha-\beta=\alpha+(-\beta)$ et
$\alpha\div\beta=\alpha\times\beta^{-1}$ pour $\beta\neq0_\Real$.
Si $\beta>0_\Real$, on définit
\[
\beta^{-1}=\Setabs{q\in\Rat}{q\leq0\ \lor\
(\exists r\in\Rat\setminus\beta)(r>0\land q<1/r)}.
\]
Si $\beta<0_\Real$, on pose $\beta^{-1}=-((-\beta)^{-1})$.
Les opposés et les produits sont ceux définis précédemment ; vérifier
que cette construction donne bien l'inverse fait partie de l'exercice.
\end{editorial}

\begin{prob}
Démontrer que $\Real$ est un corps ordonné.
\end{prob}

Il reste enfin un petit point à régler. Dans la \olref[cuts]{sec},
nous proposons de prendre
$\sqrt{2}=\Setabs{p \in \Rat}{p < 0 \text{ ou }p^2 < 2}$.
Il faut cependant montrer que cet ensemble est une \emph{coupure}.
En voici une démonstration :

\begin{proof}
Cet ensemble est un segment initial non vide et propre de $\Rat$ :
il contient $1$ et ne contient pas $2$ ; si $x<p$ et si $p$ appartient
à l'ensemble, alors soit $x<0$, soit $0\leq x<p$ et $x^2<p^2<2$.
Il suffit donc de montrer qu'il n'a pas de maximum. Si $p\leq0$,
l'élément $1$ lui est supérieur. Il suffit ainsi de traiter le cas
d'un rationnel strictement positif $p$ tel que $p^2<2$.
Posons $q=\frac{2p+2}{p+2}$ et montrons que $p<q$ et $q^2<2$.
Puisque $p+2>0$, la première inégalité résulte du calcul suivant :
\begin{align*}
  p^2 &< 2\\
  p^2 + 2p &< 2 + 2p\\
  p(p + 2) &< 2 + 2p\\
  p &< \tfrac{2+2p}{p+2} = q
\end{align*}
Pour établir $q^2<2$, il suffit de remarquer que :
\begin{align*}
  p^2 &< 2\\
  2p^2 + 4p + 2 &< p^2 + 4p+ 4\\
  4p^2 + 8p + 4 &< 2(p^2 + 4p + 4)\\
  (2p+2)^2 &< 2(p+2)^2\\
  \tfrac{(2p+2)^2}{(p+2)^2} &< 2\\
  q^2 &< 2
\end{align*}
\end{proof}
\end{document}
